\documentclass[11pt]{article}
\usepackage[utf8]{inputenc}
\usepackage{amsmath,amssymb}
\usepackage[margin=1in]{geometry}
\usepackage{hyperref}
\emergencystretch=3em

\title{Weber's equation for the fluctuation function:
$f''+(\tfrac12-2\lambda-\tfrac{z^2}{4})f=0$}
\author{Claude, with Daniel Murfet}
\date{2026-09-06}

\begin{document}
\maketitle
\sloppy

\begin{abstract}
Unit 8 of the grammar-paper \S4 autoformalisation. The substitution $S_\lambda(a)=e^{\beta a^2/8}
f(a\sqrt{\beta/2})$ turns the fluctuation ODE (iii) into \emph{Weber's equation}
$f''(z)+(\tfrac12-2\lambda-\tfrac{z^2}{4})f(z)=0$ (\texttt{prop:fluctuation\_weber}), the standard
parabolic cylinder equation with order $\nu=-2\lambda$ --- so the Weber order depends only on the
RLCT $\lambda$, not on the inverse temperature $\beta$. Proved from the \texttt{HasDerivAt} derivative
ladder alone, no special-function library. Zero \texttt{sorry}/\texttt{axiom}.
\end{abstract}

\section{Setting}
The closed form (unit~7) identified $S_\lambda$ with a parabolic cylinder function $D_{-2\lambda}$. This
proposition is the ODE face of that identification, and unlike the closed form it needs \emph{no}
$D_\nu$ machinery: it is a pure second-derivative computation on top of the already-formalised
derivative ladder ($S'_\lambda=\beta S_{\lambda+1/2}$, $S''_\lambda=\beta^2 S_{\lambda+1}$) and the
ODE (iii) $S''_\lambda=\tfrac{a\beta}{2}S'_\lambda+\lambda\beta S_\lambda$. It is the last
non-operator item of \S4 that stands independently of $D_\nu$.

\section{The thing formalised}
{\small
\begin{verbatim}
noncomputable def weberSub (β lam z : ℝ) : ℝ :=
  Real.exp (-z ^ 2 / 4) * fluctuation β lam (z * Real.sqrt (2 / β))

theorem fluctuation_weber (β lam z : ℝ) (hβ : 0 < β) (hlam : 0 < lam) :
    deriv (deriv (fun z => weberSub β lam z)) z
      + (1 / 2 - 2 * lam - z ^ 2 / 4) * weberSub β lam z = 0
\end{verbatim}
}
We formalise the substitution solved for $f$: $f(z)=e^{-z^2/4}S_\lambda(z\sqrt{2/\beta})$ (with
$\sqrt{2/\beta}=1/\sqrt{\beta/2}$, so $z=a\sqrt{\beta/2}$).

\section{Proof strategy}
Write $k=\sqrt{2/\beta}$, so the argument is $a=zk$ and $\beta k^2=2$. The plan (GPT-6 Astra) is
``prove the first derivative as an explicit function everywhere, then differentiate that function'':
\begin{enumerate}
\item Build \texttt{HasDerivAt} for the two fluctuation compositions $a\mapsto S_\lambda(ak)$ and
$a\mapsto\beta S_{\lambda+1/2}(ak)$ by \texttt{HasDerivAt.comp} with the linear map $x\mapsto xk$ (giving
the extra factor $k$), and for the Gaussian $e^{-z^2/4}$ (derivative $-\tfrac{z}{2}e^{-z^2/4}$).
\item Product rule gives $f_1(w)=\operatorname{deriv} f\,(w)=e^{-w^2/4}\bigl(-\tfrac w2S_\lambda(wk)+k\beta
S_{\lambda+1/2}(wk)\bigr)$ as a function; \texttt{funext} of \texttt{HasDerivAt.deriv} turns this into
the \emph{function equality} $\operatorname{deriv} f=f_1$.
\item Differentiate $f_1$ once more (product$+$sum$+$const-mul rules) to get $\operatorname{deriv}(\operatorname{deriv} f)(z)=f_{2,z}$.
\item Substitute the ODE (iii) for the $\beta^2S_{\lambda+1}(zk)$ term. The whole expression factors as
$e^{-z^2/4}(\beta k^2-2)\bigl(\lambda S_\lambda(zk)+\tfrac{zk}{2}\beta S_{\lambda+1/2}(zk)\bigr)$, which
vanishes because $\beta k^2=2$.
\end{enumerate}
The three fluctuation values $S_\lambda(zk)$, $S_{\lambda+1/2}(zk)$, $S_{\lambda+1}(zk)$ stay opaque
atoms throughout; the ODE eliminates the third and \texttt{ring} exposes the common factor
$\beta k^2-2$.

\section{Roadblocks and resolutions}
Three \texttt{HasDerivAt} mechanics, each a reusable lesson:
\begin{itemize}
\item \textbf{\texttt{comp} higher-order unification splits at the wrong place.} Composing the
derivative of $a\mapsto\beta S_{\lambda+1/2}(a)$ with the linear map, \texttt{exact (\ldots).comp} with
the goal function $\lambda x.\,\beta S_{\lambda+1/2}(xk)$ let the elaborator peel the outer function as
$(\beta\cdot-)$ rather than the intended $S_{\lambda+1/2}$-composition (error: expected
\texttt{HasDerivAt (HMul.hMul $\beta$) \ldots}). Fix: state that goal with the function written as an
explicit \texttt{Function.comp} $\bigl((\lambda a.\,\beta S_{\lambda+1/2}(a))\circ(\lambda x.\,xk)\bigr)$,
so no HOU is needed.
\item \textbf{Combinator derivative values need normalising.} The product/sum/\texttt{const\_mul}
combinators produce a derivative in a fixed (un-factored) shape. \texttt{convert \ldots using 1} stops
at the \texttt{ContinuousLinearMap.smulRight} wrapper (leaving a goal \texttt{ring} cannot touch), and
\texttt{using 2} leaves a spurious \texttt{Module}-instance diamond goal
(\texttt{Semiring.toModule} vs \texttt{(NormedAlgebra.toNormedSpace $\mathbb R$).toModule}). The clean
tool is \texttt{HasDerivAt.congr\_deriv h (by ring)}, which changes only the derivative value and lets
\texttt{ring} prove the raw$=$clean equality, while \texttt{exact} reconciles the instance diamond by
defeq.
\item \texttt{simpa \ldots using (comp result)} fails on the same instance diamond (its final match is
stricter than \texttt{exact}); using \texttt{exact} on the bound composition succeeds.
\end{itemize}

\section{What was Mathlib, what was new}
Everything analytic is Mathlib: \texttt{HasDerivAt.comp}/\texttt{.mul}/\texttt{.add}/\texttt{.const\_mul}/
\texttt{.exp}, \texttt{hasDerivAt\_pow}, \texttt{HasDerivAt.congr\_deriv}, \texttt{HasDerivAt.deriv}, and
\texttt{Real.sq\_sqrt}. The new content is entirely the assembly: the explicit first- and
second-derivative functions, and the observation that the ODE substitution collapses the second-order
combination to a $(\beta k^2-2)$ multiple, which the single square-root identity $\beta k^2=2$ kills.

\section{Lessons}
For a change-of-variables ODE identity, differentiate the substituted function as an \emph{explicit
function} (via \texttt{funext}$\,\circ\,$\texttt{HasDerivAt.deriv}) rather than manipulating
\texttt{deriv} pointwise --- the second derivative is then just ``differentiate a named function
again''. And \texttt{HasDerivAt.congr\_deriv} is the right way to reconcile a combinator's raw
derivative value with a hand-chosen clean form: it sidesteps both the \texttt{smulRight} wrapper and the
\texttt{Module}-instance diamond that defeat \texttt{convert} and \texttt{simpa}.

\section{Follow-ups}
The remaining \S4 items rest on the positive-order $D_n$ (Hermite$\times$Gaussian), which Mathlib does
not have as a special function though it \emph{does} have the (probabilist) Hermite polynomials:
\texttt{eq:Dn\_hermite} ($D_n(z)=e^{-z^2/4}\mathrm{He}_n(z)$ after the physicist/probabilist bridge) and
\texttt{lem:qho} (the quantum harmonic oscillator). The operator identity $[b,b^\dagger]=1$ needs a Weyl
algebra framework.
\end{document}
